\section{Exact harmonic transfer across the moving links}
\label{n:sec-transfer}

The angular program follows eigenlines of a sphere operator.  A trace of a
harmonic function on the three-dimensional manifold need not follow any of
those lines: its continuation depends on the metric throughout the ball, and
it can acquire components in arbitrarily high angular modes.  We first
express the high component as an invariant graph over the low coordinates
for the exact Dirichlet maps.  We then calculate the one entry of the
resulting finite matrix that a crossing control changes to first order.

\subsection{The exact maps and two different freezing estimates}

Fix the integer $k$, its angular program from
Proposition~\ref{n:prop-angular-program}, and the radial family of
Proposition~\ref{n:prop-radial-realization}. Write
$L(s)=-\Delta_{H(s)}$ and use its continued real orthonormal frame
$e_i(s)$ with eigenvalues $\lambda_i(s)$, $1\le i\le p=(M+1)^2$,
where $M=\lfloor\vartheta k\rfloor$ is the top degree of the fixed program.
All operators act on the common space $L^2(S^2,\mu)$.
The protected outer gap can be written
\begin{equation}
 \lambda_{p+1}^{\mathrm{ord}}(s)-\max_{i\le p}\lambda_i(s)
 \ge g_{\mathrm{gap}}>0.
 \label{n:eq-outer-angular-gap}
\end{equation}
Let $H_z(t)$ be the controlled link in
\eqref{n:controlled-angular-input}, and set
$\delta_j=\eta(j)^{3/4}=j^{-9/16}$.
The radial proposition gives, for every $t\ge0.8J$,
\begin{align}
 &a\le b(t)\le1,\quad b(t)-a=O(t^{-1/2}),\quad
   b_t,b_{tt}=O(\eta(t)^2), \label{n:eq-transfer-radial-input}\\
 &\|H_z(t)-H(s(t))\|_{C^m}=O(\eta(t)^{3/4}),\quad
   \|\partial_tH_z(t)\|_{C^m}=O(\eta(t))
   \quad(m\ge0). \label{n:eq-transfer-slow-input}
\end{align}
The same reference notation includes the negative-phase interpolation.
All needed higher derivative bounds follow from the fixed smooth bump
types. They are uniform over the complete control box and every
sufficiently late start. The radial factor is independent of $z$.
We may thus increase $J$ below while leaving the angular program and its
Ricci buffer fixed. Constants may depend on these data, in particular on
$p$; they are never required to be uniform as $k$ grows.

For \(0<r<R\), let \(\mathcal R^z_{r,R}\) send boundary data on the sphere
of radius \(R\) to the trace at radius \(r\) of their harmonic extension
to the \emph{whole} ball \(B_{g_z}(0,R)\).  Both boundary spaces are
identified with \(L^2(S^2,\mu)\).  At integer logarithmic radii set
\[
 T_j=\mathcal R^z_{e^j,e^{j+1}},\qquad
 \mathfrak d(q)=\frac{-1+\sqrt{1+4q}}2,
\]
and introduce two frozen inward maps:
\begin{equation}
 P_j^0=\exp[-\mathfrak d(b(j)^{-2}L(s(j)))],
 \qquad
 P_j^*=\exp[-\mathfrak d(b(j)^{-2}(-\Delta_{H_z(j)}))].
 \label{n:eq-frozen-inward}
\end{equation}
The first map uses the reference link; the second includes the actual
control at time \(j\).

\begin{lemma}[Whole-ball restriction and frozen comparisons]
\label{n:lem-inward-transfer}
For every control sequence, \(\mathcal R^z_{r,R}\) is positivity
preserving, fixes constants, preserves the \(\mu\)-mean, and is a contraction
on \(L^2(\mu)\) for every \(0<r<R\).  In particular these conclusions hold
for \(T_j\).  Uniformly over all controls and all sufficiently late starts,
\begin{equation}
 \|T_j-P_j^0\|_{L^2\to L^2}=o(1)
 \quad\text{as }j\longrightarrow\infty.
 \label{n:eq-operator-freeze}
\end{equation}
For the finitely many smooth reference vectors there is the sharper bound
\begin{equation}
 \|(T_j-P_j^*)e_i(s(j+1))\|_2
 \leq C\left(\eta(j)\log\frac1{\eta(j)}+\eta(j)^2\right)
 =o(\delta_j),\qquad i\leq p.
 \label{n:eq-frame-freeze}
\end{equation}
No constant or remainder in these statements depends on the controls at
earlier radii.  The \(o(\delta_j)\) assertion is only on the displayed
finite smooth frame, not in the full operator norm.
\end{lemma}

\begin{proof}
For smooth data, the classical Dirichlet theorem gives a smooth harmonic
extension \(u\) inside the ball
\cite[Theorem~6.14, p.~107]{GilbargTrudinger}.  The fixed angular area
form makes its equation away from the origin
\[
 u_{rr}+2\frac{f'}f u_r+f^{-2}\Delta_{H_z(\log r)}u=0.
\]
Consequently, for \(\bar u(r)=\int u(r,\omega)\,d\mu(\omega)\),
\[
 (f(r)^2\bar u_r(r))'=0.
\]
Regularity in the Euclidean core and \(f(r)=r\) near zero force the constant
of integration to vanish.  The mean on every inner sphere therefore equals
the boundary mean.  The maximum principle makes harmonic restriction
positive and constant preserving
\cite[Theorems~3.1 and~3.3, pp.~32--33]{GilbargTrudinger}.
A positive constant-preserving map
satisfies the pointwise square inequality
\((\mathcal R v)^2\leq\mathcal R(v^2)\): apply positivity to
\(\mathcal R((v-c)^2)\) with \(c=\mathcal Rv\) at the point in question.
After integration and mean preservation this gives
\begin{equation}
 \|\mathcal R^z_{r,R}v\|_2^2
 \leq\int\mathcal R^z_{r,R}(v^2)\,d\mu
 =\int v^2\,d\mu.
 \label{n:eq-all-radius-contraction}
\end{equation}
Density extends the maps to \(L^2(\mu)\).  For clarity, these extended maps
still are traces of one harmonic function in the interior: approximate the
boundary value in \(L^2\) by smooth data and use
\eqref{n:eq-all-radius-contraction} on inner spheres.  The local
\(L^2\)-to-sup estimate, followed by interior Schauder estimates and their
iterates for the differentiated smooth equation, makes the extensions
Cauchy in every smooth norm on smaller compact sets
\cite[Theorem~8.17, p.~194, with \(p=2\), and Corollary~6.3,
pp.~93--94]{GilbargTrudinger}.  Their limit is harmonic, and uniqueness
follows from the same approximation.
This also shows that restriction from \(L^2\) boundary data is smooth on
every strictly inner sphere.

In logarithmic radius \(t=\log r\), the equation becomes
\begin{equation}
 u_{tt}+\gamma(t)u_t+b(t)^{-2}\Delta_{H_z(t)}u=0,
 \qquad \gamma(t)=1+2b_t(t)/b(t).
 \label{n:eq-logarithmic-harmonic}
\end{equation}
For the local estimates just used and those below, its divergence form is
\[
 \partial_t(e^t b(t)^2u_t)+e^t\Delta_{H_z(t)}u=0.
\]
In fixed angular charts the second term is divergence with respect to the
fixed density \(\mu\).  On a time interval of fixed width, division by
\(e^{t_0}\) at its center leaves uniformly elliptic coefficients with the
uniform smooth bounds in the input.  The cited local estimates therefore
have constants independent of the large center time and of the controls.
The compact core is covered by ordinary smooth coordinate charts.
We first prove the operator-norm statement.  Consider arbitrary
\(j_n\to\infty\), arbitrary allowed starts \(J_n\leq j_n\) and control sequences, and data
\(\|v_n\|_2\leq1\).  Pass to a subsequence for which \(v_n\rightharpoonup v\)
weakly and \(H(s(j_n))\to H_\infty\) smoothly.  Such a phase subsequence
exists because the reference program is compact.  Move the outer boundary
to zero by \(x=t-(j_n+1)\), and write \(u_n(x)\) for the exact whole-ball
extension.  On every fixed compact interval in \(x\leq0\), the radial and
angular input estimates make the shifted coefficients converge smoothly to
those of
\begin{equation}
 U_{xx}+U_x+a^{-2}\Delta_{H_\infty}U=0.
 \label{n:eq-frozen-cylinder}
\end{equation}
Indeed, the phase changes by \(O(\eta(j_n))\) on such an interval and the
control amplitudes vanish.  The interval on which the stated end bounds
hold extends an unbounded distance to the left in these coordinates,
since \(j_n-0.8J_n\geq0.2j_n\) when \(J_n\leq j_n\).

The sectional bounds \(\|u_n(x)\|_2\leq1\) hold for every \(x<0\).
On each compact subcylinder of \(x<0\) they give an \(L^2\) bound for the
solution.  The local \(L^2\)-to-sup estimate and interior Schauder estimates
\cite[Theorem~8.17 and Corollary~6.3]{GilbargTrudinger} give uniform
\(C^{2,\alpha}\) bounds on smaller cylinders.  Compact embedding, followed by a
diagonal subsequence, gives strong \(L^2\) convergence at each interior
sphere, in particular at \(x=-1\).  The limit \(U\) solves
\eqref{n:eq-frozen-cylinder} on \((-\infty,0)\) and has sectional norm at
most one there.

The trace of this limit at \(x=0\) must also be identified.  For a fixed
smooth \(\phi\) on \(S^2\), put \(y_n(x)=\langle u_n(x),\phi\rangle_2\).
Self-adjointness of the angular Laplacian in the fixed space \(L^2(\mu)\)
gives, on \([-1,0)\),
\[
 y_n''+\gamma_n y_n'
 =-\langle u_n,b_n^{-2}\Delta_{H_{z,n}}\phi\rangle_2.
\]
The right side and \(y_n\) are uniformly bounded, and so is \(\gamma_n\).
The mean value theorem on \([-1,-1/2]\) bounds \(y_n'\) at one point.
Integrating the first-order equation for \(y_n'\) with its bounded
integrating factor then bounds \(y_n'\) on all of \([-1,0)\).
For smooth boundary data \(y_n(0)=\langle v_n,\phi\rangle_2\); the same
identity and Lipschitz bound pass to \(L^2\) data by the preceding
approximation.  Taking the limit first in \(n\) and then in \(x\uparrow0\)
shows that \(U\) has weak boundary trace \(v\).

For an eigenvalue \(\lambda\geq0\) of \(-\Delta_{H_\infty}\), the two radial
roots in \eqref{n:eq-frozen-cylinder} are
\(\mathfrak d(a^{-2}\lambda)\) and
\(-1-\mathfrak d(a^{-2}\lambda)\).  The coefficient of the latter root
must vanish because \(U\) remains sectionally \(L^2\)-bounded as
\(x\to-\infty\).  This also excludes the root \(-1\) on constants.
The weak trace fixes the remaining coefficient of every eigenfunction, so
\[
 U(x)=\exp[x\mathfrak d(a^{-2}(-\Delta_{H_\infty}))]v,\qquad x\leq0.
\]
Thus \(T_{j_n}v_n\) converges strongly to \(U(-1)\).
Apply the same argument to the frozen half-cylinder solutions
\[
 U_n^0(x)=\exp[x\mathfrak d(b(j_n)^{-2}L(s(j_n)))]v_n.
\]
They have the same sectional bound, the same weak boundary trace limit, and
the same frozen limiting equation; hence \(P_{j_n}^0v_n=U_n^0(-1)\)
has the same strong limit.  If \eqref{n:eq-operator-freeze} failed, one
could choose such \(v_n\) nearly attaining a fixed positive lower bound for
\(\|T_{j_n}-P_{j_n}^0\|\), contradicting these two strong limits.
The sequence of controls and starts was arbitrary, which proves the stated
uniformity.

For the finite-frame estimate write \(\eta_j=\eta(j)\) and set
\[
 \ell_j=4\log(1/\eta_j)=3\log j,\qquad
 \varepsilon_j=(1+\ell_j)\eta_j.
\]
Take \(J\) late enough that \(j-\ell_j>0.8J\) for every \(j\geq J\).
It suffices to check this at \(j=J\), since \(j-3\log j\) is increasing
for \(j>3\).  The entire comparison cylinder
\([j-\ell_j,j+1]\) then lies in the controlled end, including the first step.
For \(\phi_j=e_i(s(j+1))\), let
\[
 U_j(t)=\exp[(t-j-1)\mathfrak d(b(j)^{-2}(-\Delta_{H_z(j)}))]\phi_j,
 \qquad t\leq j+1.
\]
This is the bounded-backward solution of the frozen equation with boundary
value \(\phi_j\).  The vectors \(\phi_j\) have uniform smooth norms, and
the frozen metrics range over a compact smooth family.  Applying powers of
their Laplacians in the spectral expansion bounds all required spatial and
time Sobolev norms of \(U_j\), uniformly for \(t\leq j+1\).  Elliptic
estimates on the sphere and Sobolev embedding therefore give the required
uniform pointwise derivative bounds, including \(\|U_j(t)\|_\infty\leq C\).

On \([j-\ell_j,j+1]\), speeds are comparable to \(\eta_j\).  By
\eqref{n:eq-transfer-radial-input}--\eqref{n:eq-transfer-slow-input}, the
angular coefficients vary from their values at \(j\) by
\(O((1+\ell_j)\eta_j)\) in every needed spatial norm, and
\(\gamma(t)-1=O(\eta_j^2)\).  Substituting \(U_j\) into the actual operator
in \eqref{n:eq-logarithmic-harmonic} thus leaves a residual \(R_j\) with
\(\|R_j\|_\infty\leq C\varepsilon_j\).
Let \(u_j\) be the exact whole-ball extension of \(\phi_j\), and put
\(w_j=u_j-U_j\) on this cylinder.  At \(j+1\) it vanishes.  At \(j-\ell_j\)
it is uniformly bounded in sup norm: the maximum principle bounds \(u_j\)
by \(\|\phi_j\|_\infty\) throughout its ball, independently of every
earlier control, and \(U_j\) has the bound just obtained.

For late \(J\), \(\gamma\geq3/4\).  Choose one sufficiently large constant
\(K\), and define
\[
 B_j(t)=K\varepsilon_j(j+1-t)
       +K\exp[-(t-j+\ell_j)/2].
\]
It dominates \(|w_j|\) at both ends.  Since
\[
 (\partial_t^2+\gamma\partial_t)B_j
 =-K\gamma\varepsilon_j+
   K(1/4-\gamma/2)\exp[-(t-j+\ell_j)/2]
 \leq-C\varepsilon_j,
\]
the comparison principle applied to \(B_j+w_j\) and \(B_j-w_j\)
\cite[Theorem~3.3, p.~33]{GilbargTrudinger} gives
\[
 \|w_j(j)\|_2\leq\|w_j(j)\|_\infty
 \leq C(\varepsilon_j+e^{-\ell_j/2})
 =O\left(\eta_j\log(1/\eta_j)+\eta_j^2\right).
\]
The last expression divided by \(\delta_j=\eta_j^{3/4}\) tends to zero.
This proves \eqref{n:eq-frame-freeze}.  Only the local coefficient bounds
entered the residual; the whole-ball maximum principle absorbed the entire
earlier history into the bounded left discrepancy.
\end{proof}

\subsection{An invariant graph above the finite angular cluster}

Let \(\iota_j:\mathbb R^p\to L^2(\mu)\) be the isometry
\(\iota_jc=\sum_i c_i e_i(s(j))\).  Write \(E_j=\iota_j\mathbb R^p\),
let \(\Pi_j\) be its orthogonal projection, and set \(Q_j=E_j^\perp\).
We use the continued frame to write the four blocks of \(T_j\) as
\[
 \begin{array}{ll}
 T_{LL}=\iota_j^*T_j\iota_{j+1},&
 T_{LQ}=\iota_j^*T_j|_{Q_{j+1}},\\
 T_{QL}=(I-\Pi_j)T_j\iota_{j+1},&
 T_{QQ}=(I-\Pi_j)T_j|_{Q_{j+1}} .
 \end{array}
\]
Finally put
\[
 d_i(j)=\mathfrak d(b(j)^{-2}\lambda_i(s(j))),\qquad
 \alpha_i(j)=e^{d_i(j)},\qquad
 \tau_j=\min_{i\leq p}\alpha_i(j)^{-1}.
\]
There is a start-independent lower bound
\[
 \tau_*:=
 \inf_{\substack{s\geq0,\ b\in[a,1]\\ i\leq p}}
       e^{-\mathfrak d(b^{-2}\lambda_i(s))}>0,\qquad \tau_j\geq\tau_*,
\]
because \(p\) and the periodically repeated angular family are fixed.
It may be very small for the chosen \(k\).

\begin{proposition}[Invariant spaces for the exact Dirichlet maps]
\label{n:prop-exact-graphs}
After increasing \(J\), there is a fixed \(\kappa<1\) such that, uniformly
in the controls,
\begin{equation}
 \begin{aligned}
 T_{QL}&=O(\delta_j),&
 T_{LL}&=\operatorname{diag}(\alpha_i(j)^{-1})+O(\delta_j),\\
 T_{LQ}&=o(1),&
 \|T_{QQ}\|&\leq\kappa\tau_j.
 \end{aligned}
 \label{n:eq-transfer-blocks}
\end{equation}
There is a unique family of graph maps
\(V_j:\mathbb R^p\to Q_j\), among families in one fixed sufficiently small
operator ball, such that \(T_j\) maps
\[
 \operatorname{graph}_{j+1}V_{j+1}
 =\{\iota_{j+1}c+V_{j+1}c:c\in\mathbb R^p\}
\]
isomorphically onto \(\operatorname{graph}_jV_j\).
They satisfy \(\|V_j\|\leq C\delta_j\).  For each fixed \(j\), \(V_j\) is
continuous in operator norm as a function of the complete control sequence
with its countable product topology.  The outward matrix on these spaces is
\begin{equation}
 A_j=(T_{LL}+T_{LQ}V_{j+1})^{-1}
     =\operatorname{diag}(\alpha_i(j))+O(\delta_j),
 \qquad T_{LQ}V_{j+1}=o(\delta_j).
 \label{n:eq-exact-outward}
\end{equation}
It too is continuous in the product topology.
\end{proposition}

\begin{proof}
The frame moves by \(O(\eta(j))=o(\delta_j)\) over one step.  The perturbed
first-\(p\) spectral projection \(\Pi_j^*\) for \(H_z(j)\) satisfies
\(\|\Pi_j^*-\Pi_j\|=O(\delta_j)\): apply the isolated-cluster projection
formula to \eqref{n:eq-transfer-slow-input} and
\eqref{n:eq-outer-angular-gap}.  This assertion concerns the entire cluster,
so internal double crossings cause no loss.  Since \(P_j^*\) preserves
\(\Pi_j^*\), its action on the reference frame has high component
\(O(\delta_j)\).  Smooth finite-cluster functional calculus gives its low
block as \(\operatorname{diag}(\alpha_i(j)^{-1})+O(\delta_j)\).
The finite-frame estimate \eqref{n:eq-frame-freeze}, together with frame
movement, proves the first two bounds in \eqref{n:eq-transfer-blocks}.

The other two blocks need only \eqref{n:eq-operator-freeze}.  The reference
operator \(P_j^0\) preserves \(E_j\) and \(Q_j\); replacing \(Q_{j+1}\)
by \(Q_j\) costs \(O(\eta(j))\).  Hence \(T_{LQ}=o(1)\).  The gap
\eqref{n:eq-outer-angular-gap} and strict monotonicity of \(\mathfrak d\)
give a fixed \(\kappa_{\mathrm{ref}}<1\) such that
\[
 \|P_j^0|_{Q_j}\|
 \leq \kappa_{\mathrm{ref}}\tau_j .
\]
Indeed the difference between
\(\mathfrak d(b^{-2}\lambda_{p+1}^{\mathrm{ord}})\) and the largest
low frequency has a positive minimum over the compact phase and
\(b\in[a,1]\).  Choose \(\kappa_{\mathrm{ref}}<\kappa<1\).
The absolute \(o(1)\) error in \eqref{n:eq-operator-freeze} and the frame
movement can be absorbed into \((\kappa-\kappa_{\mathrm{ref}})\tau_j\)
because \(\tau_j\geq\tau_*>0\).  This is where fixing \(p\) before \(J\)
is essential.

For \(V:\mathbb R^p\to Q_{j+1}\) in a fixed small operator ball, define
\[
 M_j(V)=T_{LL}+T_{LQ}V,\qquad
 \mathcal F_j(V)=(T_{QL}+T_{QQ}V)M_j(V)^{-1}.
\]
If \(D_j=\operatorname{diag}(\alpha_i(j)^{-1})\), the relative perturbation
\[
 \|D_j^{-1}(M_j(V)-D_j)\|
 \leq\tau_*^{-1}\bigl(O(\delta_j)+o(1)\|V\|\bigr)
\]
tends uniformly to zero.  Thus \(M_j(V)\) is invertible and
\(\|M_j(V)^{-1}\|\leq\tau_j^{-1}(1+o(1))\) throughout that ball.
In particular \(\|\mathcal F_j(0)\|\leq C\delta_j\).
The inverse-difference identity gives
\begin{align*}
 \mathcal F_j(V)-\mathcal F_j(W)
 ={}&T_{QQ}(V-W)M_j(V)^{-1}\\
 &+(T_{QL}+T_{QQ}W)
     \bigl(M_j(V)^{-1}-M_j(W)^{-1}\bigr),\\
 M_j(V)^{-1}-M_j(W)^{-1}
 ={}&-M_j(V)^{-1}T_{LQ}(V-W)M_j(W)^{-1}.
\end{align*}
The first term has Lipschitz constant at most \(\kappa+o(1)\).
The second has constant \(o(1)\), since the ball is fixed,
\(\tau_*>0\), and \(T_{LQ}=o(1)\).  Choose
\(\kappa<\kappa_0<1\), and then take \(J\) late enough that every
\(\mathcal F_j\) has Lipschitz constant at most \(\kappa_0\).
Increasing \(J\) once more makes
\(\|\mathcal F_j(0)\|\) smaller than \((1-\kappa_0)\) times the ball
radius, so every transform maps that ball into itself.

For an outer cut \(N\), start with \(V_N^{[N]}=0\) and iterate
\(V_j^{[N]}=\mathcal F_j(V_{j+1}^{[N]})\) inward.  At a fixed \(j\),
two terminal choices at \(N\) differ by at most
\(C\kappa_0^{N-j}\), uniformly in every control.  Thus \(V_j^{[N]}\)
converges uniformly as \(N\to\infty\) to \(V_j\), and
\[
 V_j=\mathcal F_j(V_{j+1}),\qquad
 \|V_j\|\leq C\sum_{m\geq0}\kappa_0^m\delta_{j+m}
 \leq C'\delta_j.
\]
The last inequality uses that \(\delta_j\) decreases.  The same contraction
proves uniqueness among graph families in the chosen ball.
The fixed graph equation says exactly that
\begin{equation}
 T_j(\iota_{j+1}c+V_{j+1}c)
   =\iota_jM_j(V_{j+1})c+V_jM_j(V_{j+1})c .
 \label{n:eq-graph-identity}
\end{equation}
Invertibility of \(M_j(V_{j+1})\) proves the graph isomorphism.
Moreover \(T_{LQ}V_{j+1}=o(1)O(\delta_{j+1})=o(\delta_j)\);
inversion of the finite low matrix yields \eqref{n:eq-exact-outward}.

We give the continuity argument because it uses less than operator-norm
continuity of \(T_j\) on all of \(L^2\).  For fixed \(\ell\), the metric on
the compact ball of radius \(e^{\ell+1}\) depends on only the finitely many
controls whose bumps meet that ball.  For fixed smooth boundary data, choose
one smooth extension into the ball.  For a fixed \(0<\alpha<1\), let
\(X=\{v\in C^{2,\alpha}(\overline B):v|_{\partial B}=0\}\).
The remaining zero-boundary problem is an equation for the Dirichlet
operator \(\mathcal L_z:X\to C^\alpha(\overline B)\), which is bijective by
the classical Dirichlet theorem
\cite[Theorem~6.14, p.~107]{GilbargTrudinger}.  Its inverse is bounded by
the bounded inverse theorem.  Converging controls make these operators
converge in norm between the indicated Hölder spaces.  A Neumann series
and the inverse-difference identity then give convergence of the solutions
and of their inner traces.  Approximation by smooth data and the contraction
bound gives
\[
 T_\ell(z_n)v\longrightarrow T_\ell(z)v
 \quad\text{in }L^2
 \quad\text{for every fixed }v\in L^2.
\]
This is strong continuity.  Suppose now that finite-rank graph maps
\(V_n:\mathbb R^p\to Q_{\ell+1}\) converge to \(V\) in operator norm.
For each coordinate basis vector \(\mathbf e\),
\begin{align*}
 \|T_\ell(z_n)V_n\mathbf e-T_\ell(z)V\mathbf e\|_2
 &\leq
   \|T_\ell(z_n)(V_n\mathbf e-V\mathbf e)\|_2
   +\|(T_\ell(z_n)-T_\ell(z))V\mathbf e\|_2\\
 &\leq\|V_n\mathbf e-V\mathbf e\|_2
   +\|(T_\ell(z_n)-T_\ell(z))V\mathbf e\|_2
 \longrightarrow0.
\end{align*}
The same statement holds on the fixed low basis vectors
\(\iota_{\ell+1}\mathbf e\).  A finite basis therefore gives operator-norm
continuity for precisely the maps from \(\mathbb R^p\) that enter
\(\mathcal F_\ell\).  Their low inverses are finite matrices with the
uniform inverse bound already proved, so inversion is continuous.
Induction shows that each finite inward graph iterate is continuous in the
finitely many controls on its finitely many balls.  Its error from \(V_j\)
is at most \(C\kappa_0^{N-j}\), uniformly over the complete control product.
Given an error tolerance, first choose \(N\) for this tail and then require
closeness of those finitely many controls.  This proves product-topology
continuity of \(V_j\).  The same finite-rank estimate and finite-matrix
inversion prove continuity of \(A_j\).  In particular the argument includes
the dependence on earlier controls and the graph's dependence on arbitrarily
late controls; it asserts no unnecessary norm continuity for \(T_j\) on
the infinite-dimensional space.
\end{proof}

\subsection{The coefficient changed by a crossing control}

At a crossing we continue to use the reference frame.  Individual
eigenvectors of the perturbed link can rotate substantially at a double
eigenvalue, so they would not provide coordinates for a uniform first-order
calculation.

\begin{proposition}[Signed off-diagonal entry]
\label{n:prop-outward-crossing}
Suppose the reference lines \(i\neq h\) have an isolated linearly split
double crossing at \(s_\nu\), with positive eigenvalue \(\lambda\).
For \(L(w)=-\Delta_{h(w)}\), let
\[
 \dot L_{\nu}(s)
 =\left.\frac{d}{d\epsilon}\right|_{\epsilon=0}
       L(w(s)+\epsilon\psi_\nu),
 \qquad
 \langle e_h(s_\nu),\dot L_\nu(s_\nu)e_i(s_\nu)\rangle_2\neq0.
\]
Write $\eta_\nu=\eta(t_\nu)$, where $s(t_\nu)=s_\nu$.
Choose the fixed bump plateau inside a sufficiently small neighborhood of
\(s_\nu\).  For every mesh point whose phase is on that plateau,
\begin{equation}
 (A_j)_{hi}
 =z_\nu\eta_\nu^{3/4}m_{hi}(s(j),b(j))+o(\delta_j),
 \label{n:eq-off-diagonal-outward}
\end{equation}
uniformly over all controls, past and future.  The continuous coefficient
\(m_{hi}\) has a fixed nonzero sign and a positive lower bound in absolute
value on this neighborhood for \(b\) near \(a\).  At the crossing, with
\(\beta=b^{-2}\) and \(\alpha=e^{\mathfrak d(\beta\lambda)}\), it is
\begin{equation}
 m_{hi}(s_\nu,b)
 =\alpha\,\beta\,\mathfrak d'(\beta\lambda)
       \langle e_h(s_\nu),\dot L_\nu(s_\nu)e_i(s_\nu)\rangle_2 .
 \label{n:eq-crossing-leading-coefficient}
\end{equation}
The reversed entry has the same nonzero coefficient at the crossing.
\end{proposition}

\begin{proof}
Fix such a mesh point and temporarily write \(\beta=b(j)^{-2}\) and
\(\epsilon=z_\nu\eta_\nu^{3/4}\).  On the plateau the only active
perturbation is \(L_\epsilon=L(w(s(j))+\epsilon\psi_\nu)\).
Put
\[
 q_\beta(x)=e^{\mathfrak d(\beta x)},\qquad
 p_\beta(x)=q_\beta(x)^{-1}.
\]
Let \(\Pi_\epsilon\) be the isolated first-\(p\) projection and
\(\Pi=\Pi_0\).  The compression to \(E=\Pi L^2\) of the frozen inward
operator is \(C_-(\epsilon)=\Pi p_\beta(L_\epsilon)\Pi|_E\).
Its part passing through the high subspace \(I-\Pi_\epsilon\) is
\(O(\epsilon^2)\): the input and output projections each have norm
\(O(\epsilon)\), while \(p_\beta(L_\epsilon)\) is a contraction.
The remaining part is a smooth finite-cluster matrix.  Differentiating it
in the reference eigenbasis gives
\[
 (C_-'(0))_{hi}
 =p_\beta[\lambda_h(s(j)),\lambda_i(s(j))]
       \langle e_h(s(j)),\dot L_\nu(s(j))e_i(s(j))\rangle_2 ,
\]
where \(q[x,y]=(q(x)-q(y))/(x-y)\) for \(x\neq y\), with the continuous
extension \(q[x,x]=q'(x)\), and the same notation is used for \(p_\beta\).
This is the derivative of the entire isolated cluster matrix, including
at a double eigenvalue.

At \(\epsilon=0\), \(C_-(0)\) is diagonal with entries
\(p_\beta(\lambda_i)\).  Differentiating its inverse and using
\[
 -q_\beta(x)q_\beta(y)p_\beta[x,y]=q_\beta[x,y]
\]
shows that the off-diagonal derivative of \(C_-(\epsilon)^{-1}\) is
\[
 q_\beta[\lambda_h(s(j)),\lambda_i(s(j))]
       \langle e_h(s(j)),\dot L_\nu(s(j))e_i(s(j))\rangle_2 .
\]
This defines \(m_{hi}(s(j),b(j))\).  At equality of the eigenvalues,
\[
 q_\beta'(\lambda)
 =e^{\mathfrak d(\beta\lambda)}
       \beta\,\mathfrak d'(\beta\lambda),
\]
which proves \eqref{n:eq-crossing-leading-coefficient}.  The factor
\(\beta\) occurs exactly once, from differentiating \(q_\beta\); there is
no extra \(\beta\) in front of its divided difference.
The divided difference is continuous across equality.  The nonzero
angular entry therefore permits a fixed neighborhood on which its product
has one sign and is bounded away from zero, uniformly for \(b\) near \(a\).
Self-adjointness of \(\dot L_\nu\) in the fixed real Hilbert space makes
the reversed first-order entry equal at the crossing.

The finite-cluster Taylor remainder is \(O(\epsilon^2)\), uniformly in
the fixed phase neighborhood.  Replacing the frame at \(j+1\) by that at
\(j\) costs \(O(\eta(j))\).  Lemma~\ref{n:lem-inward-transfer} replaces the
frozen inward map on this frame by \(T_j\) with \(o(\delta_j)\) error, and
Proposition~\ref{n:prop-exact-graphs} contributes only
\(T_{LQ}V_{j+1}=o(\delta_j)\).  The inverses involved are finite matrices
with a uniform bound depending on the fixed \(\tau_*>0\); hence these
errors remain of the same order after inversion.  A fixed phase
neighborhood has time width \(O(\eta_\nu^{-1})=o(t_\nu)\), so
\(\eta(j)/\eta_\nu\to1\) uniformly there as the start tends to infinity.
Thus \(O(\epsilon^2)=O(\delta_j^2)=o(\delta_j)\), and the stated
expansion follows.  The remainders are uniform in every other control
because the frozen calculation is local, the frame comparison is uniform
over whole-ball histories, and the graph estimate is uniform over all
future tails.
\end{proof}
